% ECONOMICS_PROBLEM: {"id": "L11-566", "title": "Markups and pass-through in multiproduct firms", "jel_code": "L11", "source_id": "OP-0078"}
% !TeX program = xelatex
\documentclass[11pt,a4paper]{article}
\usepackage[margin=23mm]{geometry}
\usepackage{fontspec}
\setmainfont{Latin Modern Roman}
\usepackage{amsmath,amssymb,mathtools}
\usepackage{xcolor,enumitem,booktabs,longtable,array,xurl,hyperref}
\definecolor{muted}{HTML}{53636C}
\hypersetup{colorlinks=true,urlcolor=blue,linkcolor=blue}
\setlength{\parindent}{0pt}
\setlength{\parskip}{5pt}
\newcommand{\R}{\mathbb R}
\newcommand{\E}{\mathbb E}
\newcommand{\N}{\mathbb N}
\newcommand{\ind}{\mathbf 1}
\newcommand{\Var}{\operatorname{Var}}
\newcommand{\Cov}{\operatorname{Cov}}
\newcommand{\supp}{\operatorname{supp}}
\DeclareMathOperator*{\argmin}{arg\,min}
\DeclareMathOperator*{\argmax}{arg\,max}
\newcommand{\entrymeta}[3]{{\sffamily\small\color{muted}\textbf{\texttt{#1}}\quad|\quad #2\par #3\par}}
\newcommand{\Problemref}[1]{\texttt{#1}}
\newcommand{\JELref}[2]{\texttt{#2} (JEL \texttt{#1})}
\begin{document}
\section*{Markups and pass-through in multiproduct firms}
\textbf{Problem ID:} \texttt{L11-566}\quad \textbf{JEL:} \texttt{L11}\par
% BEGIN ECONOMICS STATEMENT
\label{problem:OP-0078}
\label{atlas:prob:I8}

\noindent\textbf{Original identifier:} I8 (atlas item 78). \textbf{Classification:} Markup and pass-through identification research program. \textbf{Original tractability label:} Medium (editorial, not a complexity result).

\subsection*{Definitions and assumptions}

For firm $f$ and product $j$, define the product markup $\mu_{fj}=p_{fj}/MC_{fj}$ when price and marginal cost are positive. Marginal cost is the derivative of a specified total-cost function with respect to product output while holding other outputs, technology, and quality fixed. For a genuinely flexible input $v$, under cost minimization, interior adjustment, a common output scale, and no unmodeled wedges, a single-output production model yields $\mu_f=\theta_f^v/\alpha_f^v$, with physical-output elasticity $\theta_f^v=\partial\log q_f/\partial\log v$ and share of revenue $\alpha_f^v=w_vv/(p_fq_f)$, where $w_v>0$ is the input price, $v>0$ the input quantity, $p_f>0$ the output price, and $q_f>0$ physical output. The identity does not automatically recover every product markup of a multiproduct firm.

\subsection*{Formal research question}

On a common firm--product sample, identify or bound product marginal costs and compare demand-side pricing first-order conditions with production-side restrictions. Let $H_{jk}=\partial q_k/\partial p_j$, and let $\Omega_{jk}$ indicate common ownership. Under static interior multiproduct Bertrand pricing,
\[
q_j+\sum_k\Omega_{jk}(p_k-MC_k)H_{jk}=0.
\]
Determine when demand derivatives, conduct, flexible-input elasticities, and input allocation are jointly identified. For a specified exogenous input-cost shifter $z$, identify the price-response vector $dp/dz$, distinguishing changes in marginal costs from quality or product scope. Characterize discrepancies that reject maintained assumptions rather than forcing the two markup measures to coincide.

\subsection*{Interpretation and scope}

A firm-level revenue production function generally combines quantities, prices, product mix, and quality. Interpreting its coefficients as physical-output elasticities requires assumptions that must be made explicit. Shared inputs and joint production make allocation of costs to individual products especially consequential. A production-derived average markup and a demand-derived product markup may therefore target different objects even in correctly estimated models. Markups are not automatically a sufficient statistic for welfare loss: fixed costs, quality investments, demand curvature, and entry also matter. Pass-through is an equilibrium derivative of prices with respect to a cost or policy shock, not generally the markup itself. Exogenous input-price shocks may alter both cost and quality, so identifying pure cost pass-through needs exclusions or a joint model. A useful reconciliation would document product definitions, cost derivatives, input adjustability, and conduct assumptions on identical data and quantify sensitivity to alternative restrictions. Rejecting a benchmark is informative but does not identify which assumption failed without additional variation.

\subsection*{Source audit and status}

[S1] verifies the production approach; [S2] studies its firm-level macroeconomic application. The original ``Berry--Gayle (2016)'' citation could not be matched to this topic and is corrected to the plausible intended Berry--Haile review [S3], whose exact authorship is verified by the publisher. This is an attribution correction, not proof of the atlas author's intent. The reconciliation target is an editorial research extension, not a claim that markups are universally unidentified.

\subsection*{References}

\begin{enumerate}[label={[S\arabic*]},leftmargin=*]

\item Jan De Loecker and Frederic Warzynski (2012). \emph{Markups and Firm-Level Export Status}. American Economic Review 102(6), 2437--2471. \url{https://www.aeaweb.org/articles?id=10.1257/aer.102.6.2437}\par \textit{Source role:} Production-side markup identification under first-order conditions. \textit{Locator:} Abstract and article metadata.

\item Jan De Loecker, Jan Eeckhout, and Gabriel Unger (2020). \emph{The Rise of Market Power and the Macroeconomic Implications}. Quarterly Journal of Economics 135(2), 561--644. \url{https://doi.org/10.1093/qje/qjz041}\par \textit{Source role:} Firm-level markup trends under production assumptions. \textit{Locator:} Abstract and article metadata.

\item Steven Berry and Philip Haile (2016). \emph{Identification in Differentiated Products Markets}. Annual Review of Economics 8, 27--52. \url{https://doi.org/10.1146/annurev-economics-080315-015050}\par \textit{Source role:} Corrected plausible reference for erroneous Berry--Gayle attribution. \textit{Locator:} Abstract and article metadata.

\end{enumerate}

\subsection*{Common conventions: Source mathematical conventions}

Unless an entry states otherwise, agent and firm sets are finite. State and action spaces are measurable, strategies and outcome maps are measurable, probability laws are specified, and expectations exist for the targets under discussion. Infinite-horizon discount factors lie in $(0,1)$ and discounted objectives are finite. Norms, tolerances and complexity input sizes must be specified for computational comparisons. Constants or primitives that appear locally are inputs, not empirical facts asserted by this catalogue.

\textbf{Identification.} A statistical model $m\in\mathcal M$ implies an observable law $P_m$ and target $\theta(m)$. At law $P$, its identified set is
\[
\Theta(P;\mathcal M)=\{\theta(m):m\in\mathcal M,\ P_m=P\}.
\]
Point identification holds when this set is a singleton, or equivalently when $P_{m_1}=P_{m_2}$ implies $\theta(m_1)=\theta(m_2)$ for all compatible models. Sharp bounds mean bounds attained, or approached when the boundary is not attained, by compatible models. Writing this set is a definition, not a solution: the substantive task is to establish informative restrictions and derive or estimate the set. An unrestricted-confounding model may give only trivial bounds.

\textbf{Inference.} Identification, estimability and valid uncertainty quantification are separate questions. Fix $\alpha\in(0,1)$. A confidence procedure $C_n$ over a declared law class $\mathcal P$ has asymptotic uniform coverage at level $1-\alpha$ when
\[
\liminf_{n\to\infty}\inf_{P\in\mathcal P}\Pr_P\{\theta(P)\in C_n\}\geq1-\alpha.
\]
For set-identified targets the coverage event must be defined separately, for example coverage of the full identified set. If an entry uses identification-indexed models instead of uniquely indexed laws, the same coverage convention is applied over those models. Pointwise limits do not establish uniform validity, and asymptotic validity does not guarantee finite-sample precision. Sample sizes, dependence structures and weak-identification sequences are part of each application.

\textbf{Counterfactuals.} $Y_i(a)$ is a potential outcome under a specified intervention, including its timing, implementation and versions. It is not $Y_i$ conditional on voluntarily choosing $a$. A full intervention vector permits interference; shorthand scalar treatment notation does not silently impose no interference. Assignment, selection, attrition, anticipation and measurement must be specified. A claim about an instrument's local effect is not automatically a population policy effect.

\textbf{Economic equilibrium.} A counterfactual model includes best responses, constraints, feasibility, market clearing, state transitions and expectations consistent with the stated information. When several equilibria exist, either a declared selector is an input or the target is set-valued. Comparisons across policy regimes must use the stated selection convention. Reduced-form relationships are not presumed invariant after an equilibrium-changing intervention.

\textbf{Optimization and welfare.} Optimization is over the stated feasible set. If attainment is not established, $\sup$ or $\inf$ and $\varepsilon$-optimal choices replace an asserted maximizer or minimizer. For example, with value $V=\sup_{a\in\mathcal A}W(a)<\infty$, an $\varepsilon$-optimal set is $\{a\in\mathcal A:W(a)\geq V-\varepsilon\}$ for $\varepsilon>0$. Benefits, costs, risk and health or environmental outcomes can be aggregated only using declared units and conversion weights. Interpersonal utility weights, the treatment of fiscal transfers and foreign profits, discounting, and the behavioral welfare criterion are explicit inputs. A comparative-static derivative additionally requires existence and appropriate differentiability of the selected counterfactual.

\textbf{Sources.} Local labels [S1], [S2], and so forth restart in every question. Locators identify the relevant abstract, model, section or result at the level actually checked; a bibliographic or abstract-level verification is not represented as inspection of an entire proof. Working papers are distinguished from later publications and changing versions. Source summaries are paraphrased. All URL checks and editorial formulations refer to the audit date stated above.
% END ECONOMICS STATEMENT
\end{document}
